\section{Non-Boolean aggregation and matrix query languages}\label{sec:matrix}
The preservation theorem of Section~\ref{sec:preservation} was proved for
an arbitrary kernel. We now compute the support number of semiring
multiplication and apply the theorem to arity elimination.

\subsection{The support number of semiring multiplication}\label{semi:width-section}
An idempotent semiring is a tuple
$\mathbb S=(S,\vee,\otimes,0_{\mathbb S},1_{\mathbb S})$: addition is a
commutative idempotent monoid, multiplication is a monoid, $0_{\mathbb S}$
is an annihilator, and multiplication distributes over finite joins on both
sides. Assume $0_{\mathbb S}\ne1_{\mathbb S}$; multiplication need not be
commutative unless stated. The \emph{join breadth}
$\breadth(\mathbb S)$ is the least $\breadth$ such that every finite join is
the join of at most $\breadth$ of its summands, or $\infty$. An idempotent
semiring whose addition orders $S$ totally --- $a\le b$ when $a\vee b=b$ ---
has breadth one, since a finite join is then its largest summand.

Semiring matrix multiplication is precisely the aggregation
\eqref{eq:aggregation} with kernel $\otimes$; we write $P\star Q$ for
$P\agg\otimes Q$. The equality array is the usual identity matrix $I_D$. The
model includes arbitrary uniform pointwise functions, constant-valued
matrices, transpose, and shared intermediates.

For this kernel the support number is governed by the join breadth.
\begin{proposition}[Semiring responses]\label{prop:semiring-width}
For semiring multiplication, the join of Definition~\ref{def:width} being
the semiring addition $\vee$,
\[
 \breadth(\mathbb S)\le\wid\otimes\le3\breadth(\mathbb S).
\]
In particular, if $\breadth<\infty$, then every computation with $m$
products has a preserving support of at most $3\breadth m$ colors, both
under restriction and under recoloring.
\end{proposition}
\begin{proof}
Distributivity gives
\[
 \ell_J(s)=\Bigl(\bigvee_{\lambda\in J}p_\lambda\Bigr)\otimes s,
 \qquad r_J(s)=s\otimes\Bigl(\bigvee_{\lambda\in J}q_\lambda\Bigr),
\]
and evaluation at $1_{\mathbb S}$ recovers the two marginal joins. The
triple of responses is therefore equivalent to the three scalar values
\[
 \bigvee_{\lambda\in J}p_\lambda,\qquad
 \bigvee_{\lambda\in J}q_\lambda,\qquad
 \bigvee_{\lambda\in J}p_\lambda\otimes q_\lambda,
\]
and it suffices to choose at most $\breadth$ colors for each of them. For
the lower bound on $\wid\otimes$, set every $q_\lambda=1_{\mathbb S}$; then
the paired response becomes $d_J=\bigvee_{\lambda\in J}p_\lambda$, so a
support $K$ must already satisfy
$\bigvee_{\lambda\in K}p_\lambda=\bigvee_{\lambda\in\Col}p_\lambda$, and a
finite family $(p_\lambda)$ whose join is not the join of fewer than
$\breadth$ of its summands therefore needs $\breadth$ colors. The
preservation statements are those of Theorems~\ref{thm:preservation}
and~\ref{thm:transport}, and require no new closure or witness calculation.
\end{proof}
The lower inequality also shows that, for these unital kernels, a finite
support number is equivalent to finite join breadth. Neither inequality is
tight in general: Boolean conjunction has $\breadth=1$ but $\wid\land=2$, by
Lemma~\ref{lem:and-width}, whereas over max-plus, min-plus and max-min the
upper bound is attained, $\wid\otimes=3$ (Appendix~\ref{app:kernels}), so
over these semirings the factor $3\breadth$ in Theorems~\ref{thm:weighted}
and~\ref{main:matrix} is the support number itself. Under restriction the
scalar maps are always held fixed; size-specific maps are allowed in the
recoloring comparison, which takes place on a single domain.

\subsection{Profile agreement over a semiring}\label{semi:weighted-section}
A general semiring has no Boolean equivalence, but the agreement test
\eqref{eq:eta} of Section~\ref{sec:phi} is already algebraic. We read the
three codes in $\mathbb S$,
\[
 \cd1=(1_{\mathbb S},0_{\mathbb S}),\quad
 \cd0=(0_{\mathbb S},1_{\mathbb S}),\quad
 \cdb=(0_{\mathbb S},0_{\mathbb S}),
\]
and let $\eta\bigl((p,q),(p',q')\bigr)=(p\otimes p')\vee(q\otimes q')$. Then
$\eta(e,e')=1_{\mathbb S}$ if $e=e'\in\{\cd0,\cd1\}$ and $0_{\mathbb S}$ if
$e\ne e'$, while the invalid code $\cdb$ matches nothing. The codes are
pairs of input values, not an extension of the scalar state space.

For a commutative idempotent semiring, we use the $2k$ independent inputs
$P_1,Q_1,\ldots,P_k,Q_k$ and define
\begin{equation}\label{eq:weighted-query}
 F_k(x,y)=\bigvee_z\bigotimes_{i=1}^k
 \bigl(P_i(x,z)\otimes P_i(z,y)\ \vee\
       Q_i(x,z)\otimes Q_i(z,y)\bigr),
\end{equation}
whose clause is $\eta$ applied to the codes of the two edges. Over the
Boolean semiring, \eqref{eq:weighted-query} is literally $\Phi_k$. Over a
general $\mathbb S$ the pairs $(P_i,Q_i)$ remain independent, and
\eqref{eq:weighted-query} defines an algebraic query on arbitrary
semiring-valued inputs rather than a Boolean equivalence.

An annotated relation assigns a scalar to each tuple. In positive annotated
relational algebra, union adds annotations, natural join multiplies them,
projection sums over removed coordinates, and renaming changes attribute
names. Following \citet{brijder2022}, ARA(3) denotes this algebra restricted
to expressions all of whose intermediate relations have at most three
attributes. The query $F_k$ has an ARA(3) expression of size $O(k)$: for
each $i$, join copies of $P_i$ renamed to $(x,z)$ and $(z,y)$, do the same
for $Q_i$, and take their union; then join the $k$ ternary results and
project away $z$. Every intermediate has at most three attributes.

The interpretation of Section~\ref{sec:phi} and Proposition~\ref{prop:semiring-width}
now give the lower bound.
\begin{theorem}[Weighted-query lower bound]\label{thm:weighted}
Let $\mathbb S$ be a nontrivial commutative idempotent semiring of finite
join breadth $\breadth$, and let $\Comp$ be a square-matrix computation of
$F_k$ on all finite inputs in which \emph{every} aggregation gate carries
the kernel $\otimes$. Writing $m$ for the number of these gates, the
\emph{products} of $\Comp$, we have
\[
 m\ \ge\ \frac{2^k}{3\breadth}\qquad(k\ge1).
\]
Conversely, $2^k$ products suffice, with $O(k2^k)$ tree nodes or $O(2^k)$
nodes with sharing. The same lower bound holds separately at every fixed
size $n\ge n_0(2^k)$, even with size-specific scalar maps and wiring.
\end{theorem}
\begin{proof}
Use the interpretation \eqref{eq:phi-inputs} of Section~\ref{sec:phi} with
the codes read in $\mathbb S$: all $N=2^k$ subsets of $[k]$ serve as colors,
$\Col=2^{[k]}$, with one anchor $x_A$ per color; the type $\tpl{x_A}$
carries the label $A$, the type $\ctype B$ carries $B$, and every other type
carries no label. All inputs are type-constant. At $(x_A,h)$ a cell witness
succeeds exactly when its color is $A$, while an anchor witness and the hub
have an unlabelled type on one of their edges, so, for every nonempty
$J\subseteq\Col$, every $f:\Col\to\Col$ and every $A\in\Col$,
\begin{equation}\label{eq:weighted-decoder}
 F_k^{\M\res J}(x_A,h)=1_{\mathbb S}\iff A\in J,\qquad
 F_k^{\M_f}(x_A,h)=1_{\mathbb S}\iff A\in f(\Col),
\end{equation}
and the value is $0_{\mathbb S}$ otherwise. All $N$ labels are detectable.

Every aggregation gate of $\Comp$ costs at most $3\breadth$ support colors
by Proposition~\ref{prop:semiring-width}, and there are $m$ of them, so
$\wid\Comp\le3\breadth m$ and the set $K$ of
Theorem~\ref{thm:preservation} has $|K|\le3\breadth m$. If $3\breadth m<N$
then $K$ misses some color $A$, and deleting the cell $\cell A$ leaves the
output at $(x_A,h)$ unchanged, whereas $F_k$ changes there by
\eqref{eq:weighted-decoder}. For the
fixed-size claim, pad with anchors carrying no label on any incident type
and apply Proposition~\ref{prop:fixed-test}. All
intermediate values may be arbitrary scalars, not just encoded bits.

For the upper bound let
$H_A=\bigotimes_{i\in A}P_i\otimes\bigotimes_{i\notin A}Q_i$ entrywise.
Distributivity and commutativity give
\begin{equation}\label{eq:weighted-upper}
 F_k=\bigvee_{A\subseteq[k]}H_A\star H_A,
\end{equation}
which is \eqref{eq:phi-upper} with $\otimes$ for $\cap$ and $\star$ for the
composition. There are $2^k$ products and $O(k2^k)$ tree nodes; sharing
prefix products gives $O(2^k)$ nodes.
\end{proof}
The lower-bound interpretation uses only the two valid codes and the invalid
one, so it also works with noncommutative multiplication if the source
product is given its displayed order; commutativity is used only for the
upper bound on general scalar inputs.

The hypothesis that every aggregation gate carries the kernel $\otimes$ is
not merely cosmetic. A computation in the sense of
Section~\ref{sec:responses} may use a different kernel at each gate, and the
support number of a computation is the sum over \emph{all} of its
aggregation gates (Definition~\ref{def:width}); a gate carrying some other
kernel therefore contributes its own support number to this sum, and a
kernel of infinite support number defeats the count altogether.
Appendix~\ref{app:kernels} exhibits such a kernel, together with a single
aggregation that uses it to compute $\Phi_k$ outright. What the
theorem counts is the number of products in a computation whose only
aggregations are products, and this is exactly the form needed for
Theorem~\ref{main:matrix}: by Lemma~\ref{lem:broadcast}, a MATLANG
computation translates into such a computation, its products contracting
the dimension $1$ having become free entrywise gates.

The following examples all have breadth one.
\begin{center}
\renewcommand{\arraystretch}{1.2}
\begin{tabular}{lcccc}
\toprule
Semiring & Scalars & $\vee$ & $\otimes$ & $(0_{\mathbb S},1_{\mathbb S})$\\
\midrule
Max--plus & $\mathbb R\cup\{-\infty\}$ & $\max$ & $+$ & $(-\infty,0)$\\
Min--plus & $\mathbb R\cup\{+\infty\}$ & $\min$ & $+$ & $(+\infty,0)$\\
Max--min & $[0,1]$ & $\max$ & $\min$ & $(0,1)$\\
\bottomrule
\end{tabular}
\end{center}
For instance the min-plus query is
\[
 \min_z\sum_{i=1}^k
 \min\{P_i(x,z)+P_i(z,y),\ Q_i(x,z)+Q_i(z,y)\},
\]
and the minimum number of min-plus products needed to compute it lies
between $2^k/3$ and $2^k$. These are not ordinary arithmetic matrix
products; the latter, combined with arbitrary decoding functions, are not
covered by this theorem.

\subsection{Typed MATLANG and arity elimination}\label{semi:typed-section}
Semiring MATLANG \citep{brijder2022}, starting from square inputs of type
$n\times n$, has intermediate types $n\times n$, $n\times1$, $1\times n$,
and $1\times1$. Besides entrywise operations and transpose it has
ones-vectors, column diagonalization, and conformable matrix
multiplication. We allow arbitrary uniform entrywise functions.

Typed MATLANG is not literally the model of Section~\ref{sec:responses}, but
it embeds into that model at no cost.
\begin{lemma}[Broadcasting typed intermediates]\label{lem:broadcast}
A typed MATLANG computation can be represented in the square-matrix model
without increasing its number of products, preserving sharing. Products
contracting dimension $n$ become one square product, and those contracting
dimension $1$ become entrywise operations.
\end{lemma}
\begin{proof}
Encode matrices, columns, rows, and scalars, respectively, as
\[
 \widehat A(u,v)=A(u,v),\quad \widehat p(u,v)=p(u),\quad
 \widehat w(u,v)=w(v),\quad \widehat s(u,v)=s,
\]
where the letters $A,p,w,s$ are local to this proof.
The scalar encoding stores the computed scalar at every entry; it is not an
input-independent free constant. Transpose and pointwise operations respect
the encodings. A ones-vector becomes the constant matrix $1_{\mathbb S}$.
Diagonalization of a column multiplies its encoding entrywise by $I_D$;
diagonalization of a scalar leaves its encoding unchanged. A product
contracting $n$ sums over the common coordinate and is one square semiring
product. A product contracting $1$ has one summand at each entry and is
entrywise $\otimes$. These identities cover all conformable type
combinations and keep the factor order. Translating each node once preserves
sharing.
\end{proof}

ARA(3) allows annotated relations of intermediate arity at most three,
whereas MATLANG uses matrices, vectors, and scalars. The following is the
precise form of Theorem~\ref{main:matrix-informal}.
\begin{theorem}[Exponential cost of arity elimination]\label{main:matrix}
Fix a nontrivial commutative idempotent semiring of finite join breadth
$\breadth$. For every $k\ge1$ there is a binary-output positive ARA(3)
expression of size $O(k)$ such that every equivalent MATLANG computation,
even with arbitrary uniform entrywise functions and sharing, uses at least
$2^k/(3\breadth)$ matrix products that contract the dimension $n$; an
equivalent MATLANG expression with $2^k$ products exists. Equivalence is on
all finite inputs over the fixed semiring.
\end{theorem}
\begin{proof}
Use the ARA(3) expression for $F_k$ and apply Lemma~\ref{lem:broadcast} to
any equivalent MATLANG computation. The translated computation has $\otimes$
at every aggregation gate, as Theorem~\ref{thm:weighted} requires, because
every product contracting the dimension $1$ has become a free entrywise
gate; that theorem then gives the lower bound, which counts only the
products contracting the dimension $n$, as stated.
Expansion~\eqref{eq:weighted-upper} is already a MATLANG expression and
gives the upper bound. The same proof applies at every fixed
$n\ge n_0(2^k)$, even when the scalar constants and entrywise maps are
chosen for that size, since they remain uniform across entries.
\end{proof}
\citet{brijder2022} establish the corresponding expressive equivalence, and
the remark following Corollary~8 in the earlier version
\citep[p.~9]{brijder2019} asks whether the exponential blow-up in arity
elimination is necessary. The theorem shows that it is, for every fixed
nontrivial commutative idempotent semiring of finite breadth, including the
three examples above, and even for the enlarged target language.

\subsection{Positive computations over further semirings}\label{semi:positive-section}
A separate transfer argument does not require idempotence. Here
\emph{positive} MATLANG uses only entrywise semiring addition and
multiplication, transpose, ones-vectors, diagonalization, and matrix
multiplication; scalar constants must lie in the subsemiring generated by
$0_{\mathbb S}$ and $1_{\mathbb S}$, and arbitrary entrywise functions are
not allowed.

On Boolean inputs such a computation is a Boolean one in disguise.
\begin{corollary}[Positive transfer]\label{cor:positive}
Let $\mathbb S$ be a nontrivial commutative semiring with
$r\cdot1_{\mathbb S}\ne0_{\mathbb S}$ for every positive integer $r$.
Interpret \eqref{eq:weighted-query} using its sums and products. Every
positive MATLANG computation of this query has at least $2^{k-1}$ products.
This holds with sharing, on all finite inputs, and separately at every fixed
size $n\ge n_0(2^k)$.
\end{corollary}
\begin{proof}
Consider the subsemiring $S_0=\{r\cdot1_{\mathbb S}:r\ge0\}$ and the map
$S_0\to\bool$, $a\mapsto[a\ne0_{\mathbb S}]$. It sends $0_{\mathbb S}$ to
$0$ and, by the assumption on $\mathbb S$, it sends $r\cdot1_{\mathbb S}$ to
$1$ for every $r\ge1$; without that assumption a positive multiple of
$1_{\mathbb S}$ could equal $0_{\mathbb S}$ and the two clauses would
conflict. The identities for integer sums and products then show that this
map is a semiring homomorphism.

On Boolean inputs all intermediate values lie in $S_0$, so applying the
homomorphism entrywise Booleanizes the computation without changing its
number of products. After Lemma~\ref{lem:broadcast} this is a Boolean
relational computation of the Boolean version of $F_k$, which is $\Phi_k$
itself. Theorem~\ref{main:phi}, or Theorem~\ref{thm:fixed-size} at a fixed
dimension, gives the bound. All transformations preserve sharing.
\end{proof}
This includes $\mathbb N,\mathbb Z,\mathbb R,\mathbb C$ with the specified
positive operations and constants, but not negative constants, subtraction,
or arbitrary decoding functions. Expansion~\eqref{eq:weighted-upper} still
uses $2^k$ products: it needs distributivity and commutativity, not
idempotence.

\paragraph{Wider scalars and other kernels.}
The invariant also measures what wider scalars and other kernels can buy.
Appendix~\ref{app:kernels} shows that $b$-bit scalars with coordinatewise
Boolean operations reduce the number of products needed for $\Phi_k$ by a
factor of $b$ and no more (Proposition~\ref{prop:width}); that every kernel
on $[0,1]$ that is monotone in each argument has support number at most
three, so the lower bounds persist with a factor $1/3$; and that a kernel of
unbounded support number, bitwise AND on the nonnegative integers, computes
$\Phi_k$ with a single aggregation. What a finite support number limits is
the way in which encoded states can interact through a witness.
